\begin{table}[htbp]
\centering
\small
\setlength{\tabcolsep}{4pt}
\sloppy
\caption{Share of measured burden borne by the most vulnerable population}
\label{tab:e3t6}
\begin{tabularx}{\textwidth}{>{\raggedright\arraybackslash}X r r r r r}
\toprule
Burden measure & Top 10\% & Top 20\% & Top 50\% & Counties & Uniform \\
\midrule
Cumulative cold exposure, employment & 19.1\% & 35.5\% & 69.6\% & 432 & no \\
Heat exposure, person-years & 14.8\% & 28.5\% & 65.4\% & 3,131 & no \\
Cold, annual Medicare cost & 14.2\% & 23.3\% & 55.9\% & 821 & no \\
Heat, annual Medicare cost & 11.1\% & 21.4\% & 52.3\% & 983 & no \\
Drought debt scar & 10.0\% & 20.0\% & 50.0\% & 603 & yes \\
2012 drought event, income & 10.0\% & 20.0\% & 50.0\% & 139 & yes \\
\bottomrule
\end{tabularx}
\begin{minipage}{\textwidth}\vspace{0.5em}\footnotesize
Counties are ordered from most to least socially vulnerable, and each column reports the share of the total measured burden borne by that fraction of the national population. A burden spread evenly across people would read 10 percent, 20 percent and 50 percent across the three columns; figures above those benchmarks mean the burden concentrates on the more vulnerable. Rows flagged as uniform per capita are built with a constant per-person burden and therefore sit on the benchmark by construction -- they are a reference line, not a finding.
\end{minipage}
\end{table}
